\section{Quadratic integral extraction from a detector}
\label{sec:quadratic-extraction}

Throughout, $A$ is a polarized abelian variety of dimension $2n$, with an integral polarization $\theta$, and an integral action of an order $\mathcal O$ in an imaginary quadratic field $K$, of signature $(n,n).$ The Rosati involution induces conjugation on $K$. Degrees mean $\deg_\theta Y=\int_Y\theta^n$ for $n$-cycles. Write $W=\bigwedge_K^{2n}H^1(A,\mathbb Q)$ for the rational Weil plane. For a nonzero $t\in W$, fix an integer $\delta\ge1$ with $\delta t\in H^{2n}(A,\mathbb Z)$, and put
\[
 T=(-1)^n\int_A t^2>0,\qquad M_T=\lceil\sqrt T\rceil.
\]
The definition of $c_t$ is the minimum of $\max(N,\deg Z^+,\deg Z^-)$ over integral effective cycles with $Nt=[Z^+]-[Z^-]$ and $N\ge1$.

We use, without reproving, ATTACK R4 (Theorem W), steps 1--3 and 5: its integral correspondence projector, orthogonality, the mass inequality, and the positive integral pairing on $W$. We also use ATTACK R1 and Theorem D for closedness of bounded complexity and bounded detection. These are supplied project inputs; their use is recorded in the check register.

\subsection{HD-D01: a quadratic refinement of Theorem W}
Choose $a\in\mathcal O$ with $a/\bar a$ not a root of unity, and use the integral projector numerator from ATTACK R4, choosing its sign so that
\[
 P([a]^*)=R e_W,\qquad R>0,\qquad P(X)=\sum_jc_jX^j\in\mathbb Z[X].
\]
Set
\[
 Q=\operatorname{Nm}(a)^n,\quad r=\operatorname{Tr}_{K/\mathbb Q}(a^{2n}),\quad
 s=4Q^2-r^2,\quad L=2Q+|r|,\quad
 m=\sum_{c_j>0}c_jQ^j.
\]
These are explicit integers, and $s>0$. Define
\[
 A_0=\left\lceil\frac{sR^2}{(n!)^2}\right\rceil,\qquad
 A_1=\left\lceil\frac{mR\sqrt{s(s+L^2)}}{n!}\right\rceil.
\]
Then the detector and division quantities satisfy the effective bound
\begin{equation}\label{eq:quadratic-general}
 \delta_W(A)\ \le\ c_t(A)\ \le\
 \delta\max\{A_0,A_1M_T\}\,\delta_W(A)^2.
\end{equation}
The upper inequality is understood on a fibre admitting a detector; the absence of a detector also excludes a finite signed representative by the supplied lower inequality.

\begin{proof}
Take an effective integral detecting $n$-cycle $Y$ of degree $D$ and put $y=P([a]^*)[Y]$. The operator $[a]^*$ has eigenvalue $Q$ on $\theta^n$, a middle summand killed by $P$, so $P(Q)=0$. Therefore
\[
 \sum_{c_j>0}c_jQ^j=\sum_{c_j<0}|c_j|Q^j=m.
\]
Pullback by $[a^j]$ preserves effectivity and multiplies the degree of an $n$-cycle by $Q^j$. Thus $y$ has integral effective positive and negative representatives of degree at most $mD$ each.

Use the positive pairing $\langle u,v\rangle=(-1)^n\int_Auv$ on $W$, and write $d=\langle y,y\rangle$. The vector $y$ is nonzero and integral, so $d\in\mathbb Z_{>0}$. The supplied mass inequality gives
\[
 d\le \frac{R^2D^2}{(n!)^2 V}\le\frac{R^2D^2}{(n!)^2},
 \qquad V=\frac1{(2n)!}\int_A\theta^{2n}\ge1.
\]
Let $A_a=[a]^*|_W$ and $B=2A_a-r$. Its two complex eigenvalues are $a^{2n},\bar a^{2n}$ before taking $2A_a-r$. The adjoint and minimal-polynomial identities are
\[
 A_a+A_a^\dagger=r,\qquad A_aA_a^\dagger=Q^2,
 \qquad B^\dagger=-B,\qquad B^2=-s.
\]
The first two identities follow from the Rosati adjoint calculation of ATTACK R4 and the two eigenlines. In particular, for $z=By$,
\[
 \langle y,z\rangle=0,\qquad\langle z,z\rangle=sd.
\]
The class $z$ is integral and has effective signed parts of degree at most $LmD$, since it is represented by $2[a]^*y-r y$.

Put $b_1=\langle\delta t,y\rangle$ and $b_2=\langle\delta t,z\rangle$. Both are integers. Orthogonality gives the exact identity
\begin{equation}\label{eq:quadratic-exact}
 (s\delta d)t=s b_1y+b_2z.
\end{equation}
This uses a common denominator $sd$ in the inverse of the diagonal Gram matrix. To estimate its cycle representatives, express $t=xu+y'v$ in the orthonormal basis $u=y/\sqrt d$, $v=z/\sqrt{sd}$. Then $x^2+(y')^2=T$, and
\[
 s|b_1|+L|b_2|
 =\delta\sqrt d\bigl(s|x|+L\sqrt s\,|y'|\bigr)
 \le\delta\sqrt d\sqrt{s(s+L^2)}\sqrt T.
\]
The two effective parts in \eqref{eq:quadratic-exact} consequently have degrees at most
\[
 \frac{\delta mR}{n!}\sqrt{s(s+L^2)}\sqrt T\,D^2.
\]
Its positive integer multiplier $s\delta d$ is at most $\delta sR^2D^2/(n!)^2$. The stated ceilings bound both quantities integrally. Minimize over detecting $Y$. The lower inequality is ATTACK R4, step 7.
\end{proof}

\subsection{HD-D02: the Gaussian fourfold constant}
For $K=\mathbb Q(i)$ with integral $\mathbb Z[i]$-action and $n=2$, every effective detector of degree $D$ gives
\begin{equation}\label{eq:gaussian-quadratic}
 c_t\le61\,440\,\delta M_T D^2.
\end{equation}
This sharper numerical form uses the following explicit correspondence, which also specifies every integral numerator.

\begin{proof}
Let $R_i=[i]^*$, $S=[1+i]^*$ and $A_a=[2+i]^*$. On the summand of $H^4$ with $j$ factors from the $i$-eigenspace, $R_i$ has eigenvalue $(-1)^j$, and $S$ has eigenvalue $(1+i)^j(1-i)^{4-j}$. Direct substitution, for $j=0,1,2,3,4$, gives
\[
 P=(1+R_i)(4-S)=16e_W.
\]
The exact substitutions are reproduced by \texttt{scripts/gaussian\_certificate.py}. For $y=P[Y]$, use
\[
 y=\operatorname{cl}(4Y+4[i]^*Y)-\operatorname{cl}([1+i]^*Y+[i(1+i)]^*Y).
\]
Each displayed effective part has degree $8D$. On $W$, the eigenvalues of $A_a$ are $-7\pm24i$. Set $z=(A_a+7)y$. It has signed parts of degree at most $(25+7)8D=256D$. With $d=\langle y,y\rangle$, the exact Gram matrix is
\[
 \operatorname{Gram}(y,z)=\begin{pmatrix}d&0\\0&576d\end{pmatrix},
 \qquad 0<d\le64D^2.
\]
For integers $b_1=\langle\delta t,y\rangle$, $b_2=\langle\delta t,z\rangle$,
\[
 (576\delta d)t=576b_1y+b_2z.
\]
In the orthonormal coordinates of $\delta t$, the signed degree bound is
\[
 D\sqrt d\sqrt{4608^2+6144^2}\,\delta\sqrt T
 =7680\delta\sqrt T D\sqrt d
 \le61\,440\delta\sqrt T D^2.
\]
The multiplier is at most $36\,864\delta D^2$. Since $M_T\ge1$, \eqref{eq:gaussian-quadratic} bounds both positive effective parts and the multiplier.
\end{proof}
